\subsection{Relatively invariant measures}

PROPOSITION 10. --- Let G be a locally compact group, \(\mu\) a relatively left-invariant measure on G with multiplier \(\chi\) . If \(\chi_{1}\) is a continuous representation of G in \(\mathbf{C}^{*}\) , the measure \(\chi_{1} \cdot \mu\) is relatively left-invariant with multiplier \(\chi_{1}\chi\) .

For,

\[
\begin{array}{r} \pmb {\gamma} (s) (\chi_ {1} \cdot \mu) = (\pmb {\gamma} (s) \chi_ {1}) \cdot (\pmb {\gamma} (s) \mu) = (\chi_ {1} (s ^ {- 1}) \chi_ {1}) \cdot (\chi (s) ^ {- 1} \mu) \\ = (\chi_ {1} \chi) (s) ^ {- 1} (\chi_ {1} \cdot \mu). \end{array}
\]

COROLLARY 1. --- Let \(\mu\) be a left Haar measure on G. For a nonzero measure \(\nu\) on G to be relatively left-invariant, it is necessary and sufficient that it be of the form \(a\chi \cdot \mu\) , where \(a \in \mathbf{C}^*\) and \(\chi\) is a continuous representation of G in \(\mathbf{C}^*\) ; its multiplier is then \(\chi\) .

The condition is sufficient (Prop. 10). On the other hand, if \(\nu\) is a nonzero relatively left-invariant measure with multiplier \(\chi\) , then \(\chi^{-1} \cdot \nu\) is left-invariant (Prop. 10) hence is of the form \(a\mu\) with \(a \in \mathbf{C}^*\) (No.~2, Cor. of Th. 1).

COROLLARY 2. --- Every relatively left-invariant measure is relatively right-invariant.

For, with the notations of Cor. 1,

\[
\begin{array}{c} \boldsymbol {\delta} (s) (\chi \cdot \mu) = \bigl (\boldsymbol {\delta} (s) \chi \bigr) \cdot \bigl (\boldsymbol {\delta} (s) \mu \bigr) = \bigl (\chi (s) \chi \bigr) \cdot \bigl (\Delta_ {\mathrm{G}} (s) \mu \bigr) \\ = (\chi \Delta_ {\mathrm{G}}) (s) (\chi \cdot \mu). \end{array}\tag{34}
\]

On account of Cor. 2, we shall henceforth speak of relatively invariant measures on G, without further specification. The relatively invariant measures admit as special cases the left Haar measures and the right Haar measures. Given a relatively invariant measure \(\nu\) on G, it is convenient to distinguish between its left multiplier \(\chi\) and its right multiplier \(\chi'\) , defined by \(\pmb{\gamma}(s)\nu = \chi(s)^{-1}\nu\) , \(\pmb{\delta}(s)\nu = \chi'(s)\nu\) . By (34), these multipliers satisfy the relation

\[
\chi^ {\prime} = \chi \Delta_ {\mathrm{G}}.\tag{35}
\]

Still denoting by \(\mu\) a left Haar measure, we have

\[
\check {\nu} = (\chi \cdot \mu) ^ {\vee} = \check {\chi} \cdot \check {\mu} = (\chi^ {- 1} \Delta_ {\mathrm{G}} ^ {- 1}) \cdot \mu ,
\]

thus \(\check{\nu}\) is relatively invariant with left multiplier \(\chi^{-1}\Delta_{\mathrm{G}}^{-1}\) and right multiplier \(\chi^{-1}\) .

The concepts of negligible, locally negligible, measurable and locally integrable function are the same for every relatively invariant measure.

